\section{Conceptual Model}
\label{sec:model}

We adopt the late-binding model documented in \path{docs/LATE_BINDING_MODEL_v1.md}.
Figure~\ref{fig:model-dag} summarizes the causal structure; Table~\ref{tab:constructs} defines core constructs.

\begin{figure}[t]
  \centering
  % Source: docs/LATE_BINDING_MODEL_v1.md (Mermaid DAG); export to PDF for camera-ready.
  \fbox{\parbox{0.92\linewidth}{\centering\vspace{2em}
    \textbf{TODO:} Include model DAG figure\\
    (\path{paper/figures/model-dag.pdf})\\
    \vspace{2em}}}
  \caption{Late-binding model: static panel truth, two-channel consumption, environmental task difficulty.}
  \label{fig:model-dag}
\end{figure}

\begin{table}[t]
  \caption{Core constructs in the late-binding model.}
  \label{tab:constructs}
  \begin{tabular}{@{}p{0.22\linewidth}p{0.70\linewidth}@{}}
\toprule
\textbf{Construct} & \textbf{Definition} \\
\midrule
Directive channel & Normative text shaping agent behavior without resolvable pointers \\
Referential channel & Pointers resolved against live repository at use time \\
Late binding & Resolution deferred to agent consumption, not authoring \\
Static truth & Mechanical VERIFIED/MISSING at pinned commit \\
Runtime resolution & Agent read/follow/act on referential pointers \\
Read--repair & Agent detects stale pointer and substitutes without external help \\
Causal load-bearing reference & Pointer on task-critical path affecting outcome \\
Environmental task difficulty & Repo-specific test/build friction gating success \\
\bottomrule
\end{tabular}

\end{table}

\paragraph{Directive channel.}
Normative text shapes agent goals and tool policy without requiring a resolvable pointer (\eg ``run tests before finishing'').
RQ5 shows this channel activates even when the referential channel is absent: mean files modified drops from $\sim$100 (A/B) to 1.7 (C) on 63 paired triplets.

\paragraph{Referential channel.}
Pointers (\paths, commands, configs) intended for runtime resolution.
Static truth is assessed mechanically in RQ1--RQ4; runtime resolution is observed in RQ5 traces (\texttt{instruction\_read}, \texttt{instruction\_followed}).

\paragraph{Late binding.}
Binding occurs at agent consumption time against the workspace snapshot, not at authoring time.
A statically false pointer may still be \emph{read} and \emph{followed} (77.8\% of condition~B runs in uptake analysis).

\paragraph{Causal load-bearing vs.\ peripheral references.}
Selection effects (85.4\% of matched pairs: cited path churn $\leq$ uncited control) suggest many cited paths are intrinsically stable; false claims may be non-load-bearing with respect to task success.
